\paragraph{Adding a source above a fixed operation.}
Fix a core operation $C$ on the nonempty passive set $W$, and let
\[
 \mathcal I(C)=\{F\text{ core}:C(v)\subseteq F(X,v)
                              \subseteq C(v)\cup X\}.
\]
The complete passive source families are fixed throughout this interval.
Only the distinguished source family $P:\mathcal D_W\to\{0,1\}$ varies.
Write $E^{\widehat k}$ for the coloring obtained from $E$ by erasing its
entire $k$-block, and put
$T_r=\{E\in\mathcal D_W:\Phi_r(C)(E_r)=1\}$.
These sets are pairwise disjoint.

\begin{theorem}[Relative source compatibility]
\label{core:thm:source-interval}
For finite $K\geq3$, each $F\in\mathcal I(C)$ is determined by its
distinguished source predicate $P$, through
\[
 F(X,v)=C(v)\cup
       \{k\in X:P(E(v)^{\widehat k})=1\}.
\]
The realizable predicates form a bounded distributive sublattice of
$\{0,1\}^{\mathcal D_W}$. Join is pointwise union, and meet is
composition at $X$:
\[
 (F\wedge G)(X,v)=F(G(X,v),v).
\]
Every realizable predicate is constant on each $T_r$.
Moreover, any predicate whose restriction to every two colorings agrees
with some member of this interval is itself realized in the interval.
\end{theorem}
\begin{proof}
At a passive-source $k$-probe, $X$ omits $k$. The two bounds therefore
force that source coefficient to equal its coefficient in $C$.
Source decomposition proves the formula and its uniqueness. At an
$X$-source probe, $C$ omits the queried winner, so $P$ is observable.
The formula shows that union takes $P\lor Q$, while composition takes
$P\land Q$. Both constructions remain in the interval. The bottom and
top are $C$ and $C\cup X$.

For constancy on $T_r$, fix $q\ne r$ and put $h(v)=F(\{q\},v)$.
This is not constantly $\{q\}$, since $C$ takes the value $\{r\}$.
The fixed-menu erasure lemma gives a nonempty core $g\leq C$.
Whenever $C=\{r\}$, necessarily $g=\{r\}$.
If $h$ ever equals $\{r\}$, preservation of its nonzero $r$-profile
therefore makes it equal $\{r\}$ everywhere on this fiber. Otherwise
it equals $\{r,q\}$ everywhere there. The $r$-profile of $F$ has the
same $X$-flag for every $q\ne r$, so this alternative is independent
of $q$. It gives the asserted constant value of $P$ on $T_r$.

Finally, let $P$ have the stated two-coloring interpolation property.
For each coloring $q$ with $P(q)=1$, intersect interval predicates
that are positive at $q$ and negative at each coloring where $P=0$.
The interpolation hypothesis supplies each factor. Their intersection
stays positive at $q$ and vanishes wherever $P=0$. Union over the
positive colorings gives exactly $P$. Empty intersections and unions
use the top and bottom.
\end{proof}

Equivalently, the interval is the downset lattice of the preorder
\[
 p\preceq_C q\quad\Longleftrightarrow\quad
 \forall F\in\mathcal I(C),\quad P_F(q)=1\ \Longrightarrow\ P_F(p)=1.
\]
The theorem reduces relative compatibility to implications between two
source probes, but this definition of the preorder still quantifies over
realizers. Antitonicity and passive cofactor independence give explicit
edges. They do not give every edge, even for a binary base.

\begin{lemma}[Three-row separation]\label{core:lem:interval-three-rows}
Fix distinct owners $a,b,c$. In variable order $X,x,y$, put
\[
 u=(\{b\},\{a\},\{c\}),\qquad
 v=(\{c\},\{b\},\{a\}),\qquad
 w=(\{a\},\{b\},\{c\}).
\]
If a core term $t$ satisfies
$t(u)\subseteq\{a,b\}$, $t(v)\subseteq\{a,c\}$, and
$t(w)\subseteq\{b,c\}$, then
$t(u)=t(v)=\{a\}$ and $t(w)=\{b,c\}$.
\end{lemma}
\begin{proof}
Suppose $b\in t(u)$. The profile pair test
\eqref{core:eq:profile-pair-safe}, applied to the bounds at $u,w$,
respectively, gives
\[
 \Phi_a(t)(\{y\})\Rightarrow\Phi_a(t)(\{x\}),\qquad
 \Phi_c(t)(\{X\})\Rightarrow\Phi_c(t)(\{y\}).
\]
The first excludes $t(v)=\{a\}$: it would force $t(u)=\{a\}$.
The second excludes $t(v)=\{c\}$: it would force $t(u)=\{c\}$.
Thus $t(v)=\{a,c\}$. Also $t(w)$ cannot be the $b$ singleton,
whose atomic coalition $\{x\}$ is the same at $v$; nor can it be the
$c$ singleton, whose coalition $\{y\}$ is the same at $u$.
Hence $t(w)=\{b,c\}$.

At the root, both $v$ and $w$ consequently act as union, so both
children respect their two bounds. At $u$, follow a child containing
$b$; if the root absorbs, it is owned by $b$ and one may follow a
$b$-singleton child. In either case the chosen child also respects
the bound at $u$. Repeat with this strictly smaller term. No variable
satisfies the three bounds and contains $b$ at $u$, a contradiction.
Thus $t(u)=\{a\}$. Swapping $x,y$ and $b,c$ gives $t(v)=\{a\}$.
The same singleton-coalition comparisons now force $t(w)=\{b,c\}$.
\end{proof}

\begin{theorem}[Closed compatibility above a binary choice]
\label{core:thm:binary-source-interval}
Let $C(x,y)=x\star_a y$ and let
$P:\mathcal D_{\{x,y\}}\to\{0,1\}$ be antitone. It is the
distinguished source predicate of an operation in $\mathcal I(C)$
if and only if:
\begin{enumerate}
\item $P$ has one value $\lambda$ on all colorings with $x=a$ or $y=a$;
\item if $\lambda=1$, then $P(x=b,y=c)=1$ for every pair of distinct
      owners $b,c\ne a$.
\end{enumerate}
\end{theorem}
\begin{proof}
The first condition follows from \cref{core:thm:source-interval}.
If the second fails, the three rows of
\cref{core:lem:interval-three-rows} would have respective outputs
$\{a,b\},\{a,c\},\{b,c\}$, which that lemma excludes.

For sufficiency with $\lambda=0$, let $Q_P$ be the greatest
single-source completion retaining $x\cup y$, given by
\cref{core:thm:source-completion}. Use $Q_P\star_a C$.
When $C=\{a\}$ the result is $C$ and the required source coefficient
is zero. Otherwise $C=x\cup y$, so $Q_P\supseteq C$ and the gate
acts as union, returning $Q_P$. This proves both interval bounds and
the required source family on every input menu.

Suppose $\lambda=1$. Every coloring with an uncolored position
extends to one with an $a$-atom, so has value one by antitonicity.
The second condition handles every remaining coloring with two
distinct foreign atoms. Thus only the values
$P(x=r,y=r)$ for $r\ne a$ can be zero.
Start from $C\cup X$ and attach $C$ by $\star_r$ once for each such
zero. This suppresses $X$ exactly when $C=\{r\}$ for a selected
foreign owner, giving precisely $P$ and preserving both bounds.
\end{proof}

For example, set $P=0$ exactly when both passive inputs are colored
by owners different from $a$. This is antitone and respects every
passive cofactor independence of $C$: fixing either input to $a$
leaves $P$ constantly one. Its singleton-$a$ value is one, but its
value at $(x=b,y=c)$ is zero. It therefore has no interval realizer.
Binary identification minors of a larger base inherit the same
additional implications.
